\chapter{Literature Survey}
\label{chap:lit}

\section{Traditional Segmentation Methods}

Before deep learning, abdominal organ segmentation leaned heavily on atlas-based techniques.
A probabilistic atlas is built from a set of manually annotated training volumes and then
deformed to fit a new scan via deformable registration. Accuracy is reasonable for large,
consistent organs like the liver, but degrades on structures with high inter-patient shape
variability --- the pancreas and adrenal glands being the canonical hard cases. Registration
itself is computationally expensive; segmenting a single volume could take tens of minutes
even on modern hardware.

Graph-cut and random-forest approaches offered a faster alternative. They exploit hand-crafted
features (Haar wavelets, HOG descriptors, local intensity histograms) and shape priors to label
voxels. The results were competitive on constrained benchmarks but required labour-intensive
feature engineering and seldom transferred well across scanner protocols or patient populations.

\section{U-Net and Encoder--Decoder Architectures}

The seminal contribution of Ronneberger et al.~\cite{ronneberger2015} was to show that an
encoder--decoder with \emph{skip connections} can learn to segment biomedical images from
relatively small annotated datasets. The encoder progressively downsamples the input, building
a hierarchy of feature maps that capture increasingly abstract representations. The decoder
mirrors this hierarchy, upsampling back to the original resolution. Skip connections concatenate
encoder feature maps directly to the corresponding decoder level, preserving fine-grained spatial
detail that would otherwise be lost through pooling.

U-Net was originally proposed for 2D histology images, but its design generalised
immediately to 3D CT and MRI. 3D U-Net variants process entire volumes through 3D convolutions,
capturing inter-slice context that 2D models miss. The cost is cubic --- doubling the
spatial extent roughly octuples the VRAM requirement --- making 3D processing impractical on
commodity GPUs for high-resolution volumes.

\section{Attention Mechanisms in Segmentation}

Attention gates, introduced by Oktay et al.~\cite{oktay2018}, address a specific weakness of
standard skip connections: they pass \emph{all} encoder features to the decoder, including
activations from background regions that add noise rather than signal. An attention gate weights
each spatial location by a learned compatibility score between the encoder feature map and a
gating signal from deeper in the network. Locations relevant to the target organ are amplified;
irrelevant regions are suppressed. The gate is differentiable and adds only a small number of
parameters.

Experimentally, attention gates improve Dice on small, irregular structures more than on large
organs --- intuitive, because large organs dominate the loss regardless, while small structures
benefit most from the noise reduction. This project's ablation confirms that pattern on the
FLARE22 dataset: adrenal gland Dice improves by approximately 0.03 when attention is enabled.

\section{nnU-Net}

Isensee et al.~\cite{isensee2021} demonstrated that a carefully configured vanilla U-Net,
with no architectural novelty, consistently outperforms bespoke methods when preprocessing,
normalisation, patch size, and post-processing are tuned systematically per dataset. Their
framework, nnU-Net, automatically configures itself from the dataset fingerprint and has topped
numerous segmentation benchmarks. It provides the natural baseline for this project: the
question is not whether a 2D model can surpass nnU-Net 3D in general, but whether it can match
it on organs where 2D context is sufficient and where the efficiency advantage of 2D processing
matters.

\section{The FLARE22 Benchmark}

The FLARE22 challenge~\cite{flare22} provides 50 CT volumes with voxel-level annotations for
13 abdominal organs. The benchmark emphasises two practical constraints: low GPU memory usage
and fast inference. Published nnU-Net 3D results on FLARE22 serve as the comparison point
throughout this report. The challenge uses a per-volume evaluation protocol --- computing Dice
per organ per volume and averaging --- which this project adopts in full.

\section{The Hausdorff Distance as a Metric}

Dice measures volumetric overlap but is insensitive to boundary quality. A model that correctly
identifies 95\% of an organ's volume but places the boundary 20~mm off in one region still
achieves a high Dice. The 95th-percentile Hausdorff Distance (HD95) captures exactly this failure
mode --- it measures the worst-case boundary error, excluding the top 5\% of outliers to
reduce sensitivity to isolated segmentation failures. Karimi and Salcudean~\cite{karimi2020}
provide a thorough analysis of HD-based loss functions and metrics in the medical imaging context.
This project reports both Dice and HD95, consistent with FLARE22 convention.

\section{Learning Rate Scheduling}

Cosine annealing, introduced by Loshchilov and Hutter~\cite{loshchilov2017}, smoothly reduces
the learning rate from its initial value to a minimum following a half-cosine curve over a fixed
number of epochs. Compared with step-based schedules, cosine annealing avoids abrupt learning
rate drops that can destabilise training, and the smooth decay has been shown empirically to
improve final accuracy across image classification and segmentation tasks. This project uses
\texttt{CosineAnnealingLR} with $T_{\max} = 60$ and $\eta_{\min} = 10^{-6}$.

\section{Summary}

The literature establishes three things relevant to this project. First, the encoder--decoder
with skip connections is a robust foundation for medical image segmentation. Second, attention
gates provide targeted improvement on small, irregular structures at low parameter cost. Third,
rigorous evaluation --- per-volume Dice, HD95, honest train/val splitting --- is as important as
architectural choice. This project combines all three insights in a practical system that runs on
freely available hardware.
